\documentclass[12pt]{article}

\usepackage{amsmath}
\usepackage{geometry}
\geometry{margin=1in}

\setlength{\parskip}{0.5em}

\begin{document}

\title{Grade 11 Physics -- Circular Dynamics Practice 2 \\ \large{Angled Forces and Vertical Circular Motion}}
\author{}
\date{}
\maketitle

\vspace{1em}

\textbf{Instructions:}

\begin{enumerate}
\item Draw a Free-Body Diagram (FBD) for every problem. If you need to break a force into components, show $F_x$ and $F_y$ clearly.
\item Write the Newton's Second Law equations ($\Sigma F_x = m a_x$, $\Sigma F_y = m a_y$) for each scenario.
\item Include units in every step and in the final answer. Use $g = 9.8 \, \text{m/s}^2$. 
\end{enumerate}

\noindent\rule{\textwidth}{0.4pt}

\vspace{1em}

\newpage
\section*{Question 1: The Conical Pendulum}

A tetherball of mass $m = 0.600 \, \text{kg}$ is attached to a rope of length $L = 2.50 \, \text{m}$. It swings in a horizontal circle such that the rope makes a constant angle of $\theta = 35.0^\circ$ with the vertical pole. 

\begin{itemize}
\item[(a)] Draw a Free-Body Diagram for the tetherball. Label the vertical and horizontal components of the tension ($T$).
\vspace{4cm}

\item[(b)] Determine the radius $r$ of the horizontal circular path.
\vspace{4cm}

\item[(c)] Calculate the tension $T$ in the rope (Hint: analyze the vertical forces first).
\vspace{5cm}

\item[(d)] Based on the horizontal net force (which acts as the centripetal force), find the speed $v$ of the tetherball.
\vspace{6cm}

\end{itemize}

\newpage

\section*{Question 2: Banked Curve (No Friction)}

A racetrack designer is planning a circular curve with a radius of $120.0 \, \text{m}$. To ensure safety in icy weather (assuming friction between tires and the icy road is exactly zero), the designer decides to bank (tilt) the curve relative to the horizontal. 

\begin{itemize}
\item[(a)] Draw the Free-Body Diagram for a car of mass $m$ on the banked curve. Which component of which force provides the inward centripetal acceleration?
\vspace{7cm}

\item[(b)] The curve is designed for cars to navigate at exactly $25.0 \, \text{m/s}$ (about $90 \, \text{km/h}$) without relying on friction. Calculate the required banking angle $\theta$. 
\textit{(Hint: You do not need the mass of the car to find this).}
\vspace{8cm}

\end{itemize}

\newpage

\section*{Question 3: Over the Hill and Through the Woods}

A $1500 \, \text{kg}$ car is driving on a hilly road. It passes over the top of a circular hill (radius = $45 \, \text{m}$) and then travels through the bottom of a circular dip (radius = $65 \, \text{m}$). 

\begin{itemize}
\item[(a)] Draw the FBD for the car exactly at the \textbf{top of the hill}. If the car is moving at $15 \, \text{m/s}$, calculate the Normal Force ($F_N$) exerted by the road on the car. Does the driver feel heavier or lighter than usual?
\vspace{6cm}

\item[(b)] Calculate the maximum speed the car can have at the top of the hill without catching air (losing contact with the road).
\vspace{5cm}

\item[(c)] The car then reaches the \textbf{bottom of the dip} while moving at $22 \, \text{m/s}$. Draw the FBD and calculate the Normal Force acting on the car.
\vspace{6cm}

\end{itemize}

\newpage

\section*{Question 4: Rollercoaster Loop}

A rollercoaster car with its passengers has a total mass of $550 \, \text{kg}$. It travels along a track that contains a vertical circular loop with a radius of $12.0 \, \text{m}$.

\begin{itemize}
\item[(a)] Draw the Free-Body Diagram for the car when it is exactly \textbf{upside-down at the top} of the loop.
\vspace{5cm}

\item[(b)] Calculate the minimum ("critical") speed the car must have at the top of the loop so that it does not fall off the track (assuming it is not locked onto the rails).
\vspace{6cm}

\item[(c)] If the rollercoaster actually travels through the top of the loop with a speed of $18 \, \text{m/s}$, calculate the normal force pushing \textit{down} from the tracks.
\vspace{6cm}

\end{itemize}

\newpage

\section*{Question 5: Artificial Gravity (Challenge)}

A cylindrical space station (radius $R = 80 \, \text{m}$) is rotating in space far from any planets to create "artificial gravity" for the astronauts inside. The astronauts stand on the inside of the outer wall, so their feet point outward from the center.

\begin{itemize}
\item[(a)] Which physical force provides the centripetal force for an astronaut? 
\vspace{4cm}

\item[(b)] To make the astronauts feel an apparent gravity equal to Earth's normal gravity ($9.8 \, \text{m/s}^2$), what must the centripetal acceleration $a_c$ of the outer wall be?
\vspace{5cm}

\item[(c)] Determine the speed $v$ at which the outer edge of the space station must rotate to simulate Earth gravity. How long, in seconds, does it take for the station to complete one full rotation?
\vspace{8cm}

\end{itemize}

\end{document}
